\subsection{Sterke comparatoren en afzonderlijke interventies}
Het primaire experiment onderscheidt A0, A1 en B1. A0 biedt technische masterdata en gangbare documentatie; A1 biedt dezelfde inhoud als B1, inclusief bronverwijzingen, tijdinformatie, gezagsbesluiten en beschikbare kwaliteitsinformatie in een conventioneel contextregister. B1 voegt de expliciete getypeerde contextbinding en machinevolgbare afhankelijkheden toe. A1 mag gestructureerde velden, SQL, zoekfuncties en rapportages gebruiken; het wordt niet kunstmatig verzwakt. Als A1 hetzelfde contract eenvoudig kan realiseren, is dat een relevante uitkomst tegen de noodzaak van B1.

De inhoud wordt vóór uitvoering vergeleken met een informatie-inventaris. Taakrelevante feiten die alleen B1 bezit worden óf aan A1 toegevoegd óf als apart pakketverschil geanalyseerd. Toegang, presentatie, zoekmogelijkheden, antwoordbudget en training worden zo gelijk mogelijk gehouden. Waar gelijkstelling niet mogelijk is, wordt het overblijvende verschil onderdeel van de interventiedefinitie. A0--B1 meet hooguit pakketwaarde; A1--B1 is de primaire toets van aanvullende contextbinding. Er wordt geen zuiver ``effect van RDF'' geclaimd wanneer tegelijk interface, tooling of verwerking verandert.

Een voorafgaande manipulatiecheck specificeert welke contextverbindingen B1 automatisch volg- en versieerbaar maakt en hoe A1 dezelfde informatie aanbiedt. Als de twee implementaties een gelijkwaardig contract en dezelfde bewerkingsmogelijkheden bieden, ontbreekt een onderscheidende semantische interventie. Het onderzoek mag dan alleen implementatiekosten/-prestaties vergelijken; een vermeend uniek SMDP-mechanisme is niet geïdentificeerd. Deze uitkomst moet vóór een causale effectclaim worden gerapporteerd.

De aanvullende ontologievariant B2 wordt afzonderlijk onderzocht. Epistemische annotatievorm en annotatiebetrouwbaarheid worden in een aparte episode gemanipuleerd. Daardoor kan een voordeel aan een specifieker mechanisme worden gerelateerd in plaats van aan alle standaarden tegelijk. De evaluatie volgt de DSR-onderscheiding tussen formatieve aanpassing en summatieve toets; de definitieve vergelijkingen worden vóór de laatste fase vastgezet. \citep{L03}

\subsection{Experimenten, variabelen en estimands}
\begin{longtable}{@{}P{.08\linewidth}P{.27\linewidth}P{.30\linewidth}P{.27\linewidth}@{}}
\caption{Geplande experimenten. Geen condities zijn uitgevoerd of resultaten verzameld.}\label{tab:experiments}\\
\toprule ID & Onafhankelijke variabele & Primaire afhankelijke variabele & Controle en grens \\ \midrule\endfirsthead
\toprule ID & Onafhankelijke variabele & Primaire afhankelijke variabele & Controle en grens \\ \midrule\endhead
E1 & Representatie A0/A1/B1; taaktype identiteit/tijd/gezag. & Aandeel inhoudelijk correcte én geldig herleidbare antwoorden; verschil B1--A1. & Gelijke informatie, toegang en training; tijd en kosten apart; gestratificeerde taken. \\
E2 & Annotatievorm conventioneel/expliciet, gekruist met juiste, ontbrekende of misleidende statusinformatie. & Onterechte gezagsopwaarderingen per beoordeelde uitspraak; correcte en onnodige onthouding. & Dezelfde bewijsinhoud per kwaliteitsniveau; mens en AI apart; geen werkelijke incidentclaim. \\
E3 & Modelreview zonder/met gerichte ontologische verdieping. & Extra juist gedetecteerde identiteit-/rolfouten én analistentijd. & Gelijke taakset en expertise; fout-positieven apart; geen labels met het correcte antwoord in de taak. \\
E4 & Zelfde versiearchief zonder/met getypeerde afhankelijkheden; intacte/verstoorde contextresolutie. & Correcte historische antwoorden en wijzigingsimpact; kritieke fouten en hersteltijd. & Dezelfde bronhistorie; verouderde mapping en ontbrekende toegang als negatieve condities. \\
\bottomrule\end{longtable}

Voor E1 is $\Delta_Q=\mathbb{E}[Q\mid B1]-\mathbb{E}[Q\mid A1]$ de beoogde gemiddelde taakuitkomst onder het vastgestelde toewijzingsontwerp. $Q=1$ vereist inhoudelijk juiste interpretatie plus een bron die het antwoord daadwerkelijk draagt; anders is $Q=0$. Componenten correctheid en herleidbaarheid worden ook afzonderlijk gerapporteerd. De gekozen representatie is de interventie; het aantal ingevulde metadatarelaties is geen consumentuitkomst. Voor E2 is $\Delta_E=\Pr(\text{onterechte opwaardering}\mid A1)-\Pr(\text{onterechte opwaardering}\mid B1)$. E3 rapporteert detectie en inspanning als vector, zonder verborgen gewogen samengestelde score. E4 scheidt impactdetectie van reconstructie en verwerkingstijd.

Het veronderstelde mechanisme is: contextbinding verandert bereikbaarheid/selectie van relevante context, die vervolgens interpretatie kan veranderen. Interfacegemak, beschikbare inhoud, modelkwaliteit, expertise en extra ontwikkelinspanning zijn rivaliserende verklaringen. Het meetontwerp registreert deze factoren; een waargenomen tussenvariabele is zonder aanvullend identificatieontwerp geen bewezen causale mediator.

\subsection{Hypothesen, praktische grenzen en weerlegging}
TP1--TP4 in het slothoofdstuk formuleren de theoretische kandidaten. De primaire operationele hypothese H1 luidt voor de vooraf afgebakende taakpopulatie: $\Delta_Q>\delta_Q$, met een toename in kritieke foutkans hoogstens $\eta$ en extra beheerlast hoogstens $\kappa$. $\delta_Q$ is de minimaal relevante winst, $\eta$ de maximaal aanvaardbare fouttoename en $\kappa$ het aanvaardbare kostenverschil. Deze grenzen worden vóór summatieve dataverzameling door verantwoordelijken vastgesteld en met rationale geregistreerd. Zolang dat niet is gebeurd, is het onderzoek niet gereed voor een bevestigende effecttoets.

H2 specificeert voor E3 een minimaal relevante detectiewinst bij een expliciete tijdsgrens. H3 specificeert voor E2 een minimale afname van onterechte opwaardering bij betrouwbare annotatie; de misleidingsconditie toetst de grens van dat voordeel. H4 betreft een vooraf relevante winst in reconstructie of impactanalyse in E4. De keuze van de primaire E4-uitkomst wordt vastgezet vóór analyse, zodat achteraf wisselen niet is toegestaan.

Ondersteuning voor de operationele centrale these T vergt een onzekerheidsinterval voor $\Delta_Q$ boven $\delta_Q$ én het halen van de kosten- en foutgrenzen. Een interval met bovengrens onder $\delta_Q$ weerspreekt de praktisch relevante winstclaim voor deze taakpopulatie; een negatief effect is sterker tegenbewijs. Praktische gelijkwaardigheid van A1 en B1 bij hogere B1-kosten weerlegt de veronderstelde noodzaak/nettowaarde van de rijkere representatie. Een breed interval dat zowel relevante winst als afwezigheid daarvan toelaat is onbeslist, niet automatisch falsificatie. Secundaire AI- of ontologiewinst redt een gefaalde primaire H1 niet achteraf.

Een nulgrens voor onterechte authenticiteitsclaims in vooraf gekozen kritieke scenario’s blijft een contractvoorstel; nul waargenomen fouten bewijst geen populatiefoutkans nul. De volledige waarden, selectie, taakclassificatie en analysemethode worden gepreregistreerd. Hier worden geen effectgroottes, poweruitkomsten of drempelwaarden verzonnen.
